\section{Introduction}
\label{sec:introduction}
Traffic cameras can support vehicle monitoring without installing a separate detector for every lane. A useful system must identify vehicle locations in individual frames, maintain identities through motion and partial occlusion, and express any speed estimate in physical units. These tasks are related, but a bounding box or a track identity alone does not establish how many metres the vehicle has travelled.

Object detectors offer complementary accuracy and computational trade-offs. Region-based detectors use proposals and region features, while one-stage detectors predict objects more directly. Transformer-based detection uses a set of object queries and assignment-based training. These families provide useful comparisons on traffic footage, particularly for distant vehicles and crowded scenes \cite{rcnn,fast,faster,rtdetr}. In a tracking-by-detection system, the detector's missed and low-confidence observations also affect trajectory continuity \cite{bytetrack}.

We present TrafficVision, a course project that connects a selection of detectors to interchangeable trackers and a camera-specific measurement pipeline. Its principal components are: (1) a UA-DETRAC annotation parser and sequence-level split; (2) custom detectors and a pretrained YOLOv8n baseline; (3) tracking and polygon zones; and (4) projection onto a road plane followed by speed estimation. The implementation, configurations and available experiment artifacts are published in the project repository on Hugging Face \cite{repository}.

This report distinguishes implemented capabilities from measured outcomes. We analyse the available \emph{validation} detection results and describe the downstream tracking and speed pipeline without treating the existence of an implementation as a quantitative validation. Reproduction of the complete experiment, a held-out test comparison and physical speed verification require the additional evidence listed in Section~\ref{sec:limitations}.
